% Focused framework presentation.
\section[Framework]{States, observables, and probability}

\begin{frame}[t]{States, observables, and probability laws}
\hypertarget{main:framework}{}
\heading{A common language for dynamics and statistics}
\begin{equation*}
  x\in V,\qquad f\in\mathcal A=B_b(V),\qquad \mu\in\mathcal P(V).
\end{equation*}
The state specifies one configuration. An observable assigns a measured value
to each state. A probability law describes an ensemble of configurations.
\heading{The pairing connects measurements to ensembles}
\begin{equation*}
  \langle f,\mu\rangle=\int_V f(x)\mu(\dd x).
\end{equation*}
\heading{The framework question}
Given an evolution of states, construct compatible evolutions of observables
and probability laws, then identify the structures they preserve.
\medskip
The examples below use linear maps, composition, and duality to carry out this construction.
\slidecite{Manuscript \S\S\,1.1--1.2, state--observable--probability framework.}
\end{frame}

\begin{frame}[t]{The algebraic structures of the framework}
\hypertarget{main:states}{}\hypertarget{main:algebras}{}
\begin{columns}[T,onlytextwidth]
  \begin{column}{.47\textwidth}
    \heading{Linear maps on states}
    \begin{equation*}
      V=\mathbb R^d,\qquad \operatorname{End}(V)\cong M_d(\mathbb R).
    \end{equation*}
    \begin{equation*}
      AB=A\circ B,\qquad I_Vx=x.
    \end{equation*}
    Composition is associative and generally noncommutative.
    \heading{Linear equations on trajectories}
    Kernels describe homogeneous solutions. A particular solution translates a kernel into an affine family.
  \end{column}
  \begin{column}{.48\textwidth}
    \heading{The observable algebra}
    Bounded Borel functions form $\mathcal A=B_b(V)$, with pointwise operations:
    \begin{equation*}
      (fg)(x)=f(x)g(x),\qquad 1(x)=1.
    \end{equation*}
    This algebra is commutative and unital.
    \heading{Functions and measures}
    The pairing $\langle f,\mu\rangle$ transfers an operator on observables to an evolution of laws.
  \end{column}
\end{columns}
\medskip
Each structure supplies a specific operation for solving or interpreting the dynamics.
\slidecite{Manuscript \S\,1.1 and \S\,1.2.1. \hyperlink{app:algebras}{Algebra operations: Appendix A1.}}
\end{frame}

\begin{frame}[t]{From a velocity field to an ODE}
\hypertarget{main:ode}{}
\fontsize{9}{11}\selectfont
A fluid element is carried by the local fluid velocity. Its position $x_t\in V$
evolves from an initial position $x$ at time $s$ according to
\begin{equation*}
  \begin{cases}
    \dot x_t=v_t(x_t),&t\in I,\\
    x_s=x.
  \end{cases}
\end{equation*}
Here $I\subseteq\mathbb R$ is an open time interval, with $(s,x)\in I\times V$.
\heading{Velocity field}
\begin{equation*}
  \begin{aligned}
    v &: I\times V\longrightarrow V,\qquad (t,z)\longmapsto v_t(z),\\
    v_t &: V\longrightarrow V,\qquad z\longmapsto v_t(z).
  \end{aligned}
\end{equation*}
At a fixed time $t$, $v_t$ assigns a velocity to every position $z$.
\heading{Trajectory}
\begin{equation*}
  x_{\cdot}:I\longrightarrow V,\qquad t\longmapsto x_t,
  \qquad x_{\cdot}\in C^1(I,V).
\end{equation*}
The solution traces one fluid element through time. Its tangent $\dot x_t$
equals the field evaluated at the current position, $v_t(x_t)$.
\slidecite{Manuscript \S\,1.2.1, fluid-motion background and ODE initial-value problems.}
\end{frame}

\begin{frame}[t]{One state flow induces compatible evolutions}
\hypertarget{main:induced}{}
\heading{State evolution}
\begin{equation*}
  x_t=\Phi_{s,t}(x_s),\qquad \Phi_{r,t}\circ\Phi_{s,r}=\Phi_{s,t}.
\end{equation*}
\heading{Pullback of observables and pushforward of laws}
\begin{equation*}
  \mathcal K_{s,t}f=f\circ\Phi_{s,t},\qquad
  \mu_t=\mathcal K_{s,t}^*\mu_s:=(\Phi_{s,t})_\#\mu_s.
\end{equation*}
\begin{equation*}
  \boxed{\langle\mathcal K_{s,t}f,\mu_s\rangle=\langle f,\mathcal K_{s,t}^*\mu_s\rangle.}
\end{equation*}
\heading{The algebra respects composition}
\begin{equation*}
  \mathcal K_{s,r}\mathcal K_{r,t}=\mathcal K_{s,t},\qquad
  \mathcal K_{r,t}^*\mathcal K_{s,r}^*=\mathcal K_{s,t}^*.
\end{equation*}
This compatibility lets us pass between trajectories, measurements, and distributions.
\slidecite{Manuscript \S\,1.1, transformation duality, and \S\,1.2.2, dynamical evolution.}
\end{frame}
